% ============================================================
%  L04_core.tex  --  Fintech Security and Regulation
%  Core Slides (10 frames)
%  Self-contained (no \input{} commands)
%  Compile: pdflatex L04_core.tex  (twice for overlays)
% ============================================================

\documentclass[aspectratio=169, 11pt]{beamer}
\usetheme{Madrid}
\usecolortheme{whale}
\usepackage{tikz,pgfplots,booktabs,multicol,amsmath,graphicx}
\usetikzlibrary{arrows.meta,positioning,shapes.geometric,calc,decorations.pathmorphing}
\pgfplotsset{compat=1.18}

% ---- Colour palette ----------------------------------------
\definecolor{mlpurple}{HTML}{9467BD}
\definecolor{mlblue}{HTML}{1F77B4}
\definecolor{mlred}{HTML}{D62728}
\definecolor{mlorange}{HTML}{FF7F0E}
\definecolor{mlgreen}{HTML}{2CA02C}
\definecolor{mlgray}{HTML}{7F7F7F}
\definecolor{mlteal}{HTML}{0D7377}
\definecolor{mlcyan}{HTML}{14BDEB}

% ---- Beamer colour settings --------------------------------
\setbeamercolor{structure}{fg=mlteal}
\setbeamercolor{palette primary}{bg=mlteal,fg=white}
\setbeamercolor{palette secondary}{bg=mlteal!80,fg=white}
\setbeamercolor{palette tertiary}{bg=mlteal!60,fg=white}
\setbeamercolor{palette quaternary}{bg=mlteal!40,fg=white}
\setbeamercolor{frametitle}{bg=mlteal!10,fg=mlteal}
\setbeamercolor{frametitle right}{bg=mlteal!5}
\setbeamercolor{block title}{bg=mlteal,fg=white}
\setbeamercolor{block body}{bg=mlteal!8,fg=black}
\setbeamercolor{block title alerted}{bg=mlred,fg=white}
\setbeamercolor{block body alerted}{bg=mlred!8,fg=black}
\setbeamercolor{block title example}{bg=mlgreen,fg=white}
\setbeamercolor{block body example}{bg=mlgreen!8,fg=black}
\setbeamercolor{title}{fg=white}
\setbeamercolor{subtitle}{fg=mlcyan}
\setbeamercolor{author}{fg=white}
\setbeamercolor{institute}{fg=white}
\setbeamercolor{date}{fg=white}

% ---- Bottom-note command -----------------------------------
\newcommand{\bottomnote}[1]{%
  \vfill
  \begin{beamercolorbox}[wd=\textwidth,ht=2ex,dp=1ex]{palette primary}%
    \tiny\hspace{1em}#1%
  \end{beamercolorbox}}

% ---- Graphics path -----------------------------------------
\graphicspath{{}}

% ---- Metadata ----------------------------------------------
\title{Fintech Security and Regulation}
\subtitle{Navigating Compliance in the Digital Finance Era}
\author{Joerg Osterrieder}
\institute{University of Zurich \\ Department of Finance}
\date{Spring 2026}
\setbeamertemplate{navigation symbols}{}

% ============================================================
\begin{document}
% ============================================================

% --- Frame 1: Title Page ---
\begin{frame}{Title Page}
  \titlepage
\end{frame}

% --- Frame 2: Bridge from Lecture 3 ---
\begin{frame}{Bridge from Lecture~3}
  \begin{columns}[T]
    \begin{column}{0.55\textwidth}
      In Lecture~3 we explored the mechanics of \textbf{payment systems}:
      card networks, real-time rails, interchange economics, and regulation
      like Durbin and PSD2.

      \vspace{0.5em}
      Now we zoom out to the full \textcolor{mlpurple}{\textbf{regulatory
      landscape}} surrounding fintech:
      \begin{itemize}
        \item \textbf{Payments regulation} was just one piece --- fintech
              firms face anti-money laundering (AML), know-your-customer
              (KYC), data protection, and licensing requirements simultaneously.
        \item \textbf{Security is existential:} A single breach can destroy
              consumer trust built over years. Regulation exists to enforce
              minimum security standards.
        \item \textbf{RegTech} is fintech's answer to the compliance burden
              --- using the same technologies (AI, APIs, cloud) to manage
              regulatory obligations at scale.
      \end{itemize}

      \vspace{0.3em}
      L04 shifts from \textcolor{mlpurple}{payment system design} to
      \textcolor{mlpurple}{the regulatory and security infrastructure that
      governs all of fintech}.
    \end{column}
    \begin{column}{0.42\textwidth}
      \includegraphics[width=\textwidth]{01_regulatory_approaches/chart.pdf}
    \end{column}
  \end{columns}
  \bottomnote{L03 showed you how money moves. L04 shows you the rules that govern every movement --- and what happens when those rules are broken.}
\end{frame}

% --- Frame 3: Regulatory Approaches Comparison ---
\begin{frame}{Regulatory Approaches to Fintech}
  \begin{columns}[T]
    \begin{column}{0.50\textwidth}
      \includegraphics[width=\textwidth]{01_regulatory_approaches/chart.pdf}
    \end{column}
    \begin{column}{0.47\textwidth}
      Jurisdictions adopt fundamentally different philosophies:

      \vspace{0.4em}
      \begin{itemize}
        \item \textbf{Activity-based} (UK, Singapore) --- Regulate what a firm
              \textit{does}, not what it \textit{is}. Enables innovation but
              creates gaps when activities span categories.
        \item \textbf{Entity-based} (USA) --- Regulate based on charter type
              (bank, broker-dealer, MSB). Clear boundaries but fintechs
              struggle to fit legacy categories.
        \item \textbf{Sandbox-first} (UK FCA, MAS, Abu Dhabi) --- Allow
              controlled experimentation with reduced requirements. Time-limited,
              small-scale, supervised.
        \item \textbf{Principles-based} (EU MiCA) --- Set broad objectives
              (consumer protection, market integrity) and let firms choose
              implementation.
      \end{itemize}

      \vspace{0.3em}
      \begin{block}{No Best Answer}
        Each approach trades off \textcolor{mlpurple}{innovation speed}
        against \textcolor{mlpurple}{consumer protection}. Most jurisdictions
        are converging toward hybrid models.
      \end{block}
    \end{column}
  \end{columns}
  \bottomnote{The UK FCA launched the first regulatory sandbox in 2016. Over 60 jurisdictions have followed. Singapore, Australia, and Hong Kong are considered leaders in sandbox design.}
\end{frame}

% --- Frame 4: Money Laundering Stages ---
\begin{frame}{Money Laundering: The Three Stages}
  \begin{columns}[T]
    \begin{column}{0.50\textwidth}
      \includegraphics[width=\textwidth]{04_money_laundering_stages/chart.pdf}
    \end{column}
    \begin{column}{0.47\textwidth}
      All money laundering follows the same three-stage pattern:

      \vspace{0.4em}
      \begin{enumerate}
        \item \textcolor{mlpurple}{\textbf{Placement}} --- Introducing illicit
              cash into the financial system. Methods: structuring deposits
              below reporting thresholds (``smurfing''), cash-intensive
              businesses, prepaid cards.
        \item \textcolor{mlpurple}{\textbf{Layering}} --- Creating complex
              transaction chains to obscure the money trail. Shell companies,
              cross-border transfers, cryptocurrency mixing, trade-based
              laundering.
        \item \textcolor{mlpurple}{\textbf{Integration}} --- Re-introducing
              ``clean'' money into the legitimate economy. Real estate,
              luxury goods, business investments, loan repayment.
      \end{enumerate}

      \vspace{0.3em}
      \begin{alertblock}{Scale of the Problem}
        The UN estimates 2--5\% of global GDP (\$800B--\$2T) is laundered
        annually. Less than 1\% is intercepted.
      \end{alertblock}
    \end{column}
  \end{columns}
  \bottomnote{Fintech creates new laundering vectors (crypto mixing, peer-to-peer transfers, neobank rapid account creation) but also new detection tools (network analysis, behavioral AI).}
\end{frame}

% --- Frame 5: KYC Process Flow ---
\begin{frame}{Know Your Customer (KYC)}
  \begin{columns}[T]
    \begin{column}{0.50\textwidth}
      \includegraphics[width=\textwidth]{03_kyc_process_flow/chart.pdf}
    \end{column}
    \begin{column}{0.47\textwidth}
      KYC is the first line of defense against financial crime:

      \vspace{0.4em}
      \begin{itemize}
        \item \textbf{Customer Identification Program (CIP)} --- Verify
              identity at onboarding: government ID, biometric checks,
              address verification, liveness detection.
        \item \textbf{Customer Due Diligence (CDD)} --- Assess risk profile:
              source of funds, expected transaction patterns, PEP
              (Politically Exposed Person) screening.
        \item \textbf{Enhanced Due Diligence (EDD)} --- For high-risk
              customers: deeper investigation, senior management approval,
              ongoing monitoring with tighter thresholds.
        \item \textbf{Ongoing monitoring} --- Transaction screening against
              sanctions lists, unusual activity detection, periodic
              re-verification.
      \end{itemize}

      \vspace{0.3em}
      \begin{block}{The Fintech Tension}
        Users expect \textcolor{mlpurple}{frictionless onboarding in minutes}.
        Regulators demand \textcolor{mlpurple}{thorough identity verification}.
        Solving this tension is a core RegTech challenge.
      \end{block}
    \end{column}
  \end{columns}
  \bottomnote{Traditional bank KYC costs \$30--\$150 per customer. AI-powered eKYC (Onfido, Jumio, Sumsub) reduces this to \$1--\$5 while improving accuracy and completion rates.}
\end{frame}

% --- Frame 6: US Regulatory Patchwork ---
\begin{frame}{The US Regulatory Patchwork}
  \begin{columns}[T]
    \begin{column}{0.50\textwidth}
      \includegraphics[width=\textwidth]{05_us_regulatory_patchwork/chart.pdf}
    \end{column}
    \begin{column}{0.47\textwidth}
      The United States has no single fintech regulator --- instead, a
      fragmented system of overlapping authorities:

      \vspace{0.4em}
      \begin{itemize}
        \item \textbf{Federal:} OCC (banks), SEC (securities), CFTC
              (derivatives), FinCEN (AML/KYC), CFPB (consumer protection),
              Federal Reserve (systemic risk).
        \item \textbf{State:} 50 separate money transmitter licenses (MTLs).
              A payments startup must apply individually in each state ---
              costing \$1--5M and 12--24 months.
        \item \textbf{Gaps:} Crypto assets may fall under SEC, CFTC, state
              regulators, or none --- depending on classification.
      \end{itemize}

      \vspace{0.3em}
      \begin{alertblock}{Regulatory Arbitrage}
        Fintechs often partner with chartered banks to ``rent'' their
        regulatory status --- the bank-as-a-service model. This creates
        \textcolor{mlred}{compliance risk concentration} at a few sponsor banks.
      \end{alertblock}
    \end{column}
  \end{columns}
  \bottomnote{The OCC proposed a federal fintech charter in 2016 but withdrew under state opposition. The dual banking system (federal + state) dates to the National Bank Act of 1863.}
\end{frame}

% --- Frame 7: Global Regulatory Comparison ---
\begin{frame}{Global Regulatory Comparison}
  \begin{columns}[T]
    \begin{column}{0.50\textwidth}
      \includegraphics[width=\textwidth]{06_global_regulatory_comparison/chart.pdf}
    \end{column}
    \begin{column}{0.47\textwidth}
      Regulatory maturity varies dramatically across jurisdictions:

      \vspace{0.4em}
      \begin{itemize}
        \item \textbf{EU:} MiCA (2024) creates unified crypto-asset
              regulation. PSD2/PSD3 for payments. GDPR for data. DORA
              (2025) for digital operational resilience.
        \item \textbf{UK:} FCA sandbox model. Activity-based licensing.
              Post-Brexit divergence creating regulatory competition with EU.
        \item \textbf{Singapore:} MAS Payment Services Act (2020). Licensing
              tiers based on transaction volume. Highly regarded sandbox.
        \item \textbf{China:} Post-2020 crackdown on big tech finance (Ant
              Group IPO halt). Strict data localization. CBDC (e-CNY)
              rollout as regulatory tool.
        \item \textbf{India:} RBI-led approach. UPI as public
              infrastructure. Account aggregator framework. Progressive
              but centralized.
      \end{itemize}
    \end{column}
  \end{columns}
  \bottomnote{MiCA is the world's first comprehensive crypto-asset regulatory framework. DORA mandates ICT risk management for all EU financial entities from January 2025.}
\end{frame}

% --- Frame 8: RegTech Technology Stack ---
\begin{frame}{RegTech: Technology Stack}
  \begin{columns}[T]
    \begin{column}{0.50\textwidth}
      \includegraphics[width=\textwidth]{07_regtech_stack_architecture/chart.pdf}
    \end{column}
    \begin{column}{0.47\textwidth}
      RegTech uses fintech's own tools to solve the compliance burden:

      \vspace{0.4em}
      \begin{itemize}
        \item \textbf{Identity layer:} Biometric verification, document
              authentication, liveness detection --- enabling eKYC in
              under 60~seconds.
        \item \textbf{Transaction monitoring:} Real-time rule engines plus
              ML-based anomaly detection. Reduces false positives from
              95\%+ (legacy) to 50--70\%.
        \item \textbf{Regulatory reporting:} Automated generation of SAR,
              CTR, and prudential reports. Machine-readable regulation
              (XBRL, RegML) enables straight-through processing.
        \item \textbf{Risk scoring:} Graph analytics for network-based risk
              (shell company detection), NLP for adverse media screening,
              geospatial analysis for sanctions compliance.
      \end{itemize}

      \vspace{0.3em}
      \begin{block}{The RegTech Paradox}
        RegTech \textit{reduces} compliance cost per transaction but
        \textit{raises} the technological bar for market entry ---
        \textcolor{mlpurple}{potentially favoring incumbents over startups}.
      \end{block}
    \end{column}
  \end{columns}
  \bottomnote{Global RegTech spending exceeded \$12B in 2024 (Juniper Research). AML transaction monitoring and identity verification are the two largest segments.}
\end{frame}

% --- Frame 9: Compliance Cost Burden ---
\begin{frame}{The Compliance Cost Burden}
  \begin{columns}[T]
    \begin{column}{0.50\textwidth}
      \includegraphics[width=\textwidth]{09_compliance_cost_comparison/chart.pdf}
    \end{column}
    \begin{column}{0.47\textwidth}
      Compliance is one of the largest operating costs for financial firms:

      \vspace{0.4em}
      \begin{itemize}
        \item \textbf{Global banks} spend 5--10\% of revenue on compliance,
              totaling over \$270B annually across the industry.
        \item \textbf{Fintechs} face disproportionate burden: compliance
              costs are largely fixed, so they consume 15--25\% of revenue
              for early-stage firms vs.\ 3--5\% for large banks.
        \item \textbf{AML fines} exceeded \$5B globally in 2023. A single
              enforcement action can be existential for a fintech
              (Wirecard: fraud, not just non-compliance).
        \item \textbf{Hidden costs:} Customer friction from KYC (20--30\%
              onboarding abandonment), compliance staff (20--30\% of
              headcount at neobanks), opportunity cost of delayed launches.
      \end{itemize}

      \vspace{0.3em}
      \begin{alertblock}{The Compliance Moat}
        High compliance costs create a \textcolor{mlred}{barrier to entry}
        that protects incumbents. Paradoxically, regulation intended to
        protect consumers can reduce competition and innovation.
      \end{alertblock}
    \end{column}
  \end{columns}
  \bottomnote{Thomson Reuters estimates the average bank employs 10--15\% of staff in compliance roles. JPMorgan alone spends over \$15B annually on regulatory and compliance functions.}
\end{frame}

% --- Frame 10: Key Takeaways ---
\begin{frame}{Key Takeaways}
  \begin{columns}[T]
    \begin{column}{0.48\textwidth}
      \textcolor{mlpurple}{\textbf{Core Concepts:}}

      \vspace{0.4em}
      \begin{enumerate}
        \item \textbf{Regulation is not monolithic} --- approaches range from
              activity-based to entity-based, sandbox to prohibition.
              No single model dominates globally.
        \item \textbf{AML/KYC is the universal requirement} --- every fintech
              touching money must implement placement, layering, and
              integration detection. The cost is non-trivial.
        \item \textbf{The US patchwork is uniquely burdensome} --- 50-state
              licensing plus overlapping federal agencies creates high
              barriers and regulatory arbitrage incentives.
        \item \textbf{RegTech is essential, not optional} --- AI-powered
              compliance tools are the only way to manage obligations
              at fintech scale and speed.
        \item \textbf{Compliance is a competitive moat} --- fixed costs
              favor incumbents; the regulatory burden is the largest
              non-technical barrier to fintech entry.
      \end{enumerate}
    \end{column}
    \begin{column}{0.48\textwidth}
      \textcolor{mlpurple}{\textbf{Looking Ahead to Lecture~5:}}

      \vspace{0.4em}
      With payment mechanics (L03) and the regulatory landscape (L04)
      established, Lecture~5 will explore \textbf{blockchain and
      decentralized finance (DeFi)} --- technologies that attempt to
      bypass traditional intermediaries and regulators entirely.

      \vspace{0.5em}
      \textbf{Questions to carry forward:}
      \begin{itemize}
        \item Can decentralized protocols comply with AML/KYC requirements
              designed for centralized institutions?
        \item Does ``code is law'' replace or complement legal regulation?
        \item How should regulators treat DeFi --- ignore, adapt, or
              prohibit?
      \end{itemize}

      \vspace{0.5em}
      \begin{block}{Evaluation Framework}
        For any fintech innovation, always ask: \textcolor{mlpurple}{Who
        regulates it, what compliance costs does it bear, and what happens
        when something goes wrong?}
      \end{block}
    \end{column}
  \end{columns}
  \bottomnote{L04 completes the regulatory foundation. L05 introduces blockchain and DeFi, where many of these regulatory assumptions are challenged or circumvented entirely.}
\end{frame}

\end{document}
